\section{Đối tượng và phạm vi nghiên cứu}
Đối tượng nghiên cứu của khóa luận là hệ thống giám sát và điều khiển môi trường trong nhà dựa trên cảm biến, truyền thông không dây, nền tảng IoT và ứng dụng di động. Hệ thống được định hướng cho căn hộ chung cư, nơi số lượng node không lớn nhưng yêu cầu tính ổn định, khả năng lắp đặt thuận tiện và giao diện sử dụng rõ ràng. Trong hệ thống này, cây lan ý được xem là thành phần sinh học hỗ trợ cải thiện vi khí hậu trong điều kiện chăm sóc phù hợp; phần IoT giữ vai trò đo lường, lưu trữ dữ liệu, cảnh báo và điều khiển.

Phạm vi nghiên cứu tập trung vào một node cảm biến và một gateway. Node được bố trí gần khu vực trồng lan ý hoặc khu vực sinh hoạt cần giám sát; các thông số đo bao gồm TDS, pH, nhiệt độ, độ ẩm, CO2 và ánh sáng. Phần điều khiển gồm quạt, đèn thông qua relay và điều hòa thông qua phát lệnh hồng ngoại. Khóa luận không đi sâu vào mô hình sinh học của cơ chế hấp thụ VOC, không định lượng tốc độ hấp thụ từng chất ô nhiễm và không đặt mục tiêu thay thế thiết bị đo IAQ chuyên dụng. Thay vào đó, đề tài tập trung xây dựng một nguyên mẫu có khả năng thu thập dữ liệu, đưa ra cảnh báo và thực hiện điều khiển trong điều kiện căn hộ thực tế.
